\section{Anforderungen}
\label{chapter4}


\subsection{Anfoderungen bezüglich der Gefahrlosigkeit}




	\begin{enumerate}[label*=\arabic*.]
		\item \label{req.1}  Anfoderung: \\
		Beim Programmstart müssen alle Funktionen bezüglich der Gefahrlosigkeit getestet werden.  \\
		Folge Anfoderungen: \ref{req.1.1} , \ref{req.1.2}
		\begin{enumerate}[label*=\arabic*.]
			\item \label{req.1.1}  Anfoderung  \\
			Beim Programmstart muss die Error-LED getestet werden.\ref{req.2}\\
			\item \label{req.1.2} Anfoderung   \\
			Beim Programmstart muss das Fehlersenden vom UART getestet werden.\\
		\end{enumerate}
		\item \label{req.2} Anforderung:   \\
		\item \label{req.3} Anforderung:
	\end{enumerate}

\subsection{Sicherheitsbezogene Anforderungen}
Bei diesem Projekt tauschen 2 getrennte Systeme Daten aus. Diese Übertragung findet über Bluetooth statt, ein normalerweise unsicheres oder schwer abzusicherndes Übertragungsmedium. Folgende Sicherheits-Primitiva sollen durch diese Implementierung abgesichert werden:
	\begin{enumerate}[label*=\arabic*.]
		\item \label{req.1}  Anfoderung: Vertraulichkeit (Confidentiality)\\
		\\
		Folge Anfoderungen: \ref{req.1.1} , \ref{req.1.2}
		\begin{enumerate}[label*=\arabic*.]
			\item \label{req.1.1}  Anfoderung  \\
			Beim Programmstart müssen die beiden Systeme neues, gemeinsames Schlüsselmaterial austauschen.\\
			\item \label{req.1.2} Anfoderung   \\
			Die initiale Kommunikation sowie die nachfolgende Kommunikation müssen abgesichert sein (Handshake, TLS, hybride Kommunikation).\\
		\end{enumerate}
		\item \label{req.2} Anforderung: Integrität (Integrity)\\
		Durch die kabellose Übertragung oder durch Fremdeinwirkung kann es zur Änderung der Daten kommen, was erkannt werden muss.\\
		\begin{enumerate}[label*=\arabic*.]
			\item \label{req.2.1}  Anfoderung  \\
			Die übertragenen Pakete oder Daten müssen stets auf geänderten Inhalt (Checksums) überprüft werden.\\
		\end{enumerate}
		\item \label{req.3} Anforderung: Verfügbarkeit (Availability)\\
		Die beiden validen Systeme müssen zu jeder Zeit zur (sicheren) gemeinsamen Kommunikation fähig sein.\\
		\begin{enumerate}[label*=\arabic*.]
			\item \label{req.3.1}  Anfoderung  \\
			Die Verschlüsselung und Sicherheitsmechanismen dürfen die Leistung nicht beeinträchtigen.\\
			\item \label{req.3.2}  Anfoderung  \\
			Die Latenz darf durch Sicherheitsmechanismen nicht negativ beeinflusst werden.\\
		\end{enumerate}
		\item \label{req.4} (Optional?) Anforderung: Authentizität (Authenticity)\\
		Es muss immer klar sein, von welchem der involvierten Systeme die gesendeten Daten stammen.
		\begin{enumerate}[label*=\arabic*.]
			\item \label{req.4.1}  Anfoderung \\
			Jedes System muss eindeutig durch Schlüsselmaterial identifizierbar sein.\\
		\end{enumerate}
	\end{enumerate}

\subsection{Anfoderungen ohne Einfluss auf Sicherheit und Gefahrenlosigkeit}

